%------------------------------------------------------------------------------%
% \chapter{Revisão}
%------------------------------------------------------------------------------%

%------------------------------------------------------------------------------%
\begin{exercicios}{Logaritmos}
    
    \begin{exercicio}
        Coloque a expressão dada na forma logarítmica
        \begin{mathlist}[2]
            \item 2^6 = 64
            \item 3^5 = 243
            \item 3^{-2} = \frac{1}{9}
            \item \left(\frac{1}{3}\right)^1 = \frac{1}{3}
            \item \left(\frac{1}{2}\right)^1{-4} = 16 
            \item 32^{3/5} = 8
            \item 81^{3/4} = 27
            \item 10^{-3} = 0,0001
        \end{mathlist}
    \end{exercicio}
    
    \begin{exercicio}
        Use o fato que $\log 3 = 0,4771$ e $\log 4 = 0,6021$ para encontrar o
        valor do logaritmo solicitado
        \begin{mathlist}[2]
            \item \log 12
            \item \log \frac{3}{4}
            \item \log 16
            \item \log \sqrt{3}
        \end{mathlist}
    \end{exercicio}
        
    \begin{exercicio}
        Utilize as propriedades dos logaritmos para simplificar as expressões dadas
        \begin{mathlist}[2]
            \item \log x(x+1)^4
            \item \log\frac{\sqrt{x+1}}{x^2+1}
            \item \ln\frac{e^x}{1+e^x}
            \item \ln xe^{-x^2}
            \item \ln x(x+1)(x+2)
            \item \ln \frac{x^{1/2}}{x^2\sqrt{1+x^2}}
            \item \ln x^x
        \end{mathlist}
        \begin{answers}
            \begin{mathlist}[2]
                \item \log x + 4\log(x+1)
                \item \frac{1}{2}\log(x+1) - \log\left(x^2+1\right)
                \item x -\ln\left(1+e^x\right)
                \item \ln x -x^2
                \item \ln x + \ln(x+1) + \ln(x+2)
                \item \frac{1}{2}\ln x - 2\ln x - \frac{1}{2}\ln\left(1+x^2\right)
                \item x\ln x
            \end{mathlist}
        \end{answers}
    \end{exercicio}
        
    \begin{exercicio}
        Use logaritmos para resolver as equações em~$t$
        \begin{mathlist}[2]
            \item e^{0,4t} = 8
            \item \frac{1}{3} e^{-3t} = 0,9
            \item 5e^{-2t} = 6
            \item 4e^{t-1} = 4
            \item \frac{50}{1+4e^{0,2t}} = 20
            \item \frac{200}{1+3e^{-0,3t}} = 100
            \item A = B e^{-t/2}
            \item \frac{A}{1+Be^{t/2}} = C
        \end{mathlist}
        \begin{answers}
          \begin{mathlist}[2]
              \item t = \frac{5}{2} \ln{\left (8 \right )}
              \item t = -\frac{1}{3}\ln(2,7)
              \item t = -\frac{1}{2}\ln{\left (\frac{6}{5} \right )}
              \item t = 1 
              \item t = 5 \ln{\left (\frac{3}{8} \right )}
              \item t = \frac{10}{3}\ln(3)
              \item t = -2\ln\left(\frac{A}{B}\right)
              \item t = 2\ln\left(\frac{A-C}{CB}\right)
            \end{mathlist}
        \end{answers}        
    \end{exercicio}
    
\end{exercicios}

%------------------------------------------------------------------------------%
